\chapter*{Introduction générale}
\addcontentsline{toc}{chapter}{Introduction générale}
\markboth{Introduction générale}{}
% Mesure sur le build du 20-10-2026 : le « Plan » lu par chaque chapitre pointe
% une fente .mtc en retard d'une division — le chapitre 1 ne trouve rien, le
% chapitre 2 affiche le plan du 1, le 3 celui du 2, le 4 celui du 3.
% \adjustmtc est la commande que minitoc fournit pour avancer ce decompte ; il
% en manque un ici, la ou le modele en a deja pose un apres le sommaire.
\adjustmtc

La fonction Ressources Humaines prend une place de plus en plus stratégique dans les entreprises, portée par la digitalisation accélérée des processus métiers. Le recrutement, première étape du cycle de vie du collaborateur, conditionne directement la qualité du capital humain d’une organisation — et pourtant, dans beaucoup de structures, ce processus reste fragmenté entre des outils hétérogènes : messagerie électronique, tableurs, formulaires papier. Cette dispersion nuit à la traçabilité des décisions, à la collaboration entre intervenants, et à l’expérience vécue par les candidats.

C’est dans ce contexte que Sopra HR Software Tunisie m’a confié le projet JOB4YOU, dans le cadre de mon Projet de Fin d’Études : concevoir et développer une plateforme web centralisée, sécurisée et intuitive, couvrant l’intégralité du processus de recrutement interne, de la publication d’une offre jusqu’à la décision finale d’embauche, en passant par la gestion des candidatures, l’organisation des entretiens et la collecte des retours des managers.

Ce rapport retrace la démarche suivie tout au long du projet, en quatre chapitres. Le premier pose le cadre général — organisme d’accueil, étude de l’existant, problématique et méthodologie adoptée. Le deuxième analyse et spécifie les besoins fonctionnels et non fonctionnels de la plateforme. Le troisième détaille la conception de la solution, sur le plan architectural comme sur celui des modèles de données. Le quatrième décrit la réalisation technique : environnement de travail et principales fonctionnalités développées. Le rapport se conclut par un bilan du travail effectué et les perspectives d’évolution envisagées pour la plateforme.
